\subsection{Porównanie czasowego rozkładu alarmów siedmiu metod}
\label{sec:seven_alarm_timelines}

W celu zbadania rzeczywistego czasu ostrzegania porównano dwie reguły
statystyczne (LPPLS Confidence Indicator i BSADF-gap) oraz pięć rodzin
algorytmów uczenia maszynowego: regresję logistyczną, drzewo decyzyjne,
XGBoost, LightGBM i CatBoost. Wszystkie metody prognozowały spadek ceny
BTC o co najmniej 10\% w ciągu 60 dni, na \textbf{1401 identycznych
datach testowych} (2021--2024), z 764 dodatnimi etykietami.

Reguły LPPLS-CI i BSADF pochodzą z osobnej procedury kalibracji
statystycznej, a pięć modeli ML (każdy w wariancie Price+LPPLS)
z eksperymentu z chronologicznym strojeniem hiperparametrów.
Wyniki są zatem porównywalne czasowo, ale \textbf{nie stanowią eksperymentu
z identycznym budżetem strojenia ani wspólnym ograniczeniem czasu alarmu}.
Pokazano rzeczywiste wcześniej wybrane decyzje OOS, nie dobierając
progów ponownie do zbioru testowego.

\begin{table}[htbp]\centering\small
\caption{Out-of-sample alarm coverage for seven methods evaluated on the same BTC 10\%-decline, 60-day target and identical daily test calendar.}
\label{tab:seven_method_alarm_shares}
\begin{tabular}{@{}lrr@{}}\toprule
Metoda & Dni w alarmie & Time in Alarm (\%) \\ \midrule
LPPLS Confidence Indicator & 188 & 13,4 \\
BSADF statistical rule & 442 & 31,5 \\
Logistic Regression & 1399 & 99,9 \\
Decision Tree & 1072 & 76,5 \\
XGBoost & 1297 & 92,6 \\
LightGBM & 1323 & 94,4 \\
CatBoost & 1354 & 96,6 \\
\bottomrule\end{tabular}
\end{table}

\begin{figure}[htbp]\centering
\includegraphics[width=\linewidth]{figures/oos_seven_methods_compact.pdf}
\caption{BTC-USD close, retrospective UNION episodes, and out-of-sample
warning states of seven methods on identical dates (2021--2024).
Statistical rules and chronologically tuned ML classifiers follow different
calibration protocols; the plot is not a matched-budget ranking.}
\label{fig:seven_method_compact}
\end{figure}

Zauważalny jest silny kontrast między regułą LPPLS-CI, alarmującą
przez 13,4\% czasu, regułą BSADF, aktywną przez 31,5\% czasu,
a większością modeli ML, alarmujących przez 76,5--99,9\% czasu.
Przy niemal ciągłym alarmowaniu wysoki recall i F1 mogą mieć
ograniczone znaczenie operacyjne. Różowe obszary UNION są
retrospektywnym odniesieniem literaturowym, nie etykietą
przyszłego spadku używaną do oceny przedstawionych alarmów.

Szczegółowe panele z czarną linią ceny i odcinkami alarmów dla
każdej metody udostępniono w dodatku do manuskryptu.
